%

%%%

%%%   Самарский государственный технический университет

%%%   Кафедра прикладной математики и информатики

%%%

%%%   Преамбула для статьи в журнал "Вестник Самарского государственного

%%%   технического университета. Серия: «Физико-математические науки»"

%%%   Ориентирована под MikTEx 2.4 for Win32

%%%

%%%   Хотя сам журнал собирается под GNU/Linux на дистрибутиве TeXLive



\documentclass[11pt,a4paper,twoside]{article}



\usepackage[cp1251]{inputenc}

\usepackage[T2A]{fontenc}

\usepackage[english,russian]{babel}

\usepackage{samgtubib}

\usepackage[dvips]{graphicx}

\usepackage[multiple]{footmisc}



\usepackage{samgtu}

\renewcommand{\VestnikYear}{2009} % Год издания Вестника

\renewcommand{\VestnikNumber}{2\,(19)} % Номер Вестника

\renewcommand{\VestnikMonth}{Ноябрь} % Месяц выхода Вестника

\renewcommand{\VestnikData}{01.11.09} % Дата подписания Вестника в печать

\renewcommand{\VestnikUPL}{26,64} % Усл. печ. л.

\renewcommand{\VestnikUIL}{25,73} % Уч.-изд. л.

\renewcommand{\VestnikRegNumber}{479/09} % Регистрационный номер

\renewcommand{\VestnikOrder}{994} % Номер заказа







%%%%

%%%   Самарский государственный технический университет

%%%   Кафедра прикладной математики и информатики

%%%

%%%   Преамбула для статьи в журнал "Вестник Самарского государственного

%%%   технического университета. Серия: «Физико-математические науки»"

%%%   Ориентирована под MikTEx 2.4 for Win32

%%%

%%%   Хотя сам журнал собирается под GNU/Linux на дистрибутиве TeXLive



\documentclass[11pt,a4paper,twoside]{article}



\usepackage[cp1251]{inputenc}

\usepackage[T2A]{fontenc}

\usepackage[english,russian]{babel}

\usepackage{samgtubib}

\usepackage[dvips]{graphicx}

\usepackage[multiple]{footmisc}



\usepackage{samgtu}

\renewcommand{\VestnikYear}{2009} % Год издания Вестника

\renewcommand{\VestnikNumber}{2\,(19)} % Номер Вестника

\renewcommand{\VestnikMonth}{Ноябрь} % Месяц выхода Вестника

\renewcommand{\VestnikData}{01.11.09} % Дата подписания Вестника в печать

\renewcommand{\VestnikUPL}{26,64} % Усл. печ. л.

\renewcommand{\VestnikUIL}{25,73} % Уч.-изд. л.

\renewcommand{\VestnikRegNumber}{479/09} % Регистрационный номер

\renewcommand{\VestnikOrder}{994} % Номер заказа







\usepackage{qrcode}

%
%\newcommand{\totop}[1]{\raisebox{6pt}[1pt][2pt]{#1}}
%\newcommand{\tottop}[1]{\raisebox{12pt}[1pt][2pt]{#1}}
%\newcommand{\totttop}[1]{\raisebox{18pt}[1pt][2pt]{#1}}
%\newcommand{\tottttop}[1]{\raisebox{24pt}[1pt][2pt]{#1}}
%\newcommand{\totttttop}[1]{\raisebox{30pt}[1pt][2pt]{#1}}
\newcommand{\tobot}[1]{\raisebox{-6pt}[1pt][2pt]{#1}}
%\newcommand{\tobbot}[1]{\raisebox{-12pt}[1pt][2pt]{#1}}
%\newcommand{\tobbbot}[1]{\raisebox{-18pt}[1pt][2pt]{#1}}

%\graphicspath{{./00PIC/}}

%\samgtupubl % для генерации pdf, отдаваемого в типографию СамГТУ
\pdfcrop % обрезка pdf для размещения на math-net.ru
\draft % На корректуре
\filllastpage % Последняя страница не заполнена не полностью
%\briefcomm % Статья является кратким сообщением

\vsgtu{1508}

\FirstPage{1}
\LastPage{24}
%
%
\begin{document}


\renewcommand{\baselinestretch}{0.9}

\hfill \qrcode[hyperlink,height=.8in]{http://dx.doi.org/10.14498/vsgtu1508}
\vspace{-.8in}


% \Partition{Механика деформируемого твердого тела}{Applied mechanics} % Раздел Вестника, в который подаётся статья
% Названия разделов см. на сайте при подаче заявки


%%%%%%%%%%%%%%%%%%%%%%%%%%%%%%%%%%%%%%%%%%%%%%%%%%%%%%%%%%%%%%%%%%%%%%%%%%%%%%%%%%%%
%Информация о статье и об авторах на русском и английском языках

% Номер УДК
\udc{539.313:517.968.72}%Обработка текста. Подготовка текстов : Языки программирования. Компьютерные языки
\msc{???}% Коды MSC см. http://www.ams.org/msc/

\rutitle{Применение линейных дробных аналогов реологических моделей в задаче аппроксимации экспериментальных данных по растяжению поливинилхлоридного пластиката}% Полный заголовок
{Применение линейных дробных аналогов реологических моделей \dots}% Заголовок, который идёт в верхний колонтитул

\ruauthor{Л.\,Г.~Унгарова}% Автор(ы) для заголовка, если авторов несколько укать \inst
{У\,н\,г\,а\,р\,о\,в\,а~Л.\,Г.}% Автор(ы) для оглавления и колонтитула через \,

\ruaffil{\rusamgtu.}
\enaffil{\ensamgtu.}
% Если несколько, то так  \institute{ Организация 1 \and Организация 2  \and Организация 3}

%\email{E-mails}{petuhovds@mail.ru, kie@icmm.ru} %E-mail's авторов

\ruauthordetails{1}{\fullauthor{\rupid{??????}{Унгарова Луиза Гадильевна}} \regalia{ \mem{algluiza@gmail.com}; автор, ведущий переписку}
\position{аспирант} \dept{каф. прикладной математики и информатики}}

\enauthordetails{1}{\fullauthor{\enpid{?????}{Luiza G. Ungarova}} \regalia{\mem{algluiza@gmail.com}; Corresponding Author}
\position{Postgraduate Student}
\dept{Dept. of Applied Mathematics \& Computer Science}}


\rukeywords{дробные реологические модели, идентификация параметров, оператор Римана"--~Лиувилля, экспериментальные данные, поливинилхлоридный пластикат, погрешность}

\ruabstract{Рассматриваются и анализируются одноосные феноменологические модели вязкоупругого деформирования на основе дробных аналогов реологических моделей Скотт Блэра, Фойхта, Максвелла, Кельвина и Зенера. Получены аналитические решения соответствующих дифференциальных уравнений с дробными операторами Римана"--~Лиувилля при действии постоянных напряжений с последующей разгрузкой, которые записываются через обобщенную (двухпараметрическую) дробную экспоненциальную функцию и в зависимости от типа модели содержат от двух до четырех параметров. Разработана методика идентификации параметров моделей на основе базовой информации для экспериментальных кривых ползучести при постоянных напряжениях. Нелинейная задача параметрической идентификации решается двухступенчаым итерационным способом. На первом этапе используются характерные экспериментальные точки диаграмм и особенности поведения моделей при неограниченном возрастании времени и определяются начальные приближения параметров.
На втором этапе осуществляется уточнение этих параметров методом покоординатного спуска (методом Хука"--~Дживса) и минимизации функционала среднеквадратического отклонения расчетных значений от экспериментальных. Методика идентификации реализована для всех рассмотренных моделей на основании известных экспериментальных данных одноосного вязкоупругого деформирования поливинилхлоридного пластиката при температуре 20\,\textcelsius\ и пяти уровней растягивающего напряжения. Приводятся табличные значения параметров для всех моделей. Выполнен анализ погрешности построенных феноменологических моделей по экспериментальным данным для всего ансамбля кривых вязкоупругого деформирования. Установлено, что для модели Скотт Блэра погрешность составляет 14.17\%, Фойхта~--- 11.13\%, Максвелла~--- 13.02\%, Кельвина~--- 10.56\%, Зенера~--- 11.06\%. Приводятся графики расчетных и экспериментальных зависимостей вязкоупругой деформации  для поливинилхлоридного пластиката.
}

\entitle{The use of linear fractional analogues rheological models in the problem of approximating the experimental data on the stretch Polyvinylchloride Elastron}
%
\enauthor{L.\,G.~Ungarova, }{U\,n\,g\,a\,r\,o\,v\,a~L.\,G.}


\enabstract{Considered and analyzed the uniaxial phenomenological models of viscoelastic deformation based on fractional rheologic models analogues of Scott Blair, Voigt, Maxwell, Kelvin and Zener. Analytical solutions of  the corresponding differential equations are obtained with fractional Riemann--Liouville operators under constant stress with further unloading, that are written by the generalized (two-parameter) fractional exponential function and contains from two to four parameters depending on the type of model. A method for identifying the model parameters based on the background information for the experimental creep curves with constant stresses. Nonlinear problem parametric identification  is solved by two-step iterative method. The first stage uses the characteristic data points diagrams and features in the behavior of the models under unrestricted growth of time and the initial approximation of parameters are determined. At the second stage, the refinement of these parameters by coordinate descent (a method of Hooke--Jeeves) and minimizing the functional standard deviation for calculated and experimental values. Method of identification is realized for all the considered models on the basis of the known experimental data uniaxial viscoelastic deformation of Polyvinylchloride Elastron at a temperature of 20\,\textcelsius\ and five the tensile stress levels. Table-valued parameters for all models are given. The errors analysis constructed phenomenological models is made to experimental data over the entire ensemble of curves viscoelastic deformation. It was found that the model Scott Blair's error is 14.17\%, Voigt --- 11.13\%, Maxwell --- 13.02\%, Kelvin --- 10.56\%, Zener --- 11.06\%. The graphs of the calculated and experimental dependences viscoelastic deformation of Polyvinylchloride Elastron are submitted.}

\enkeywords{fractional analogue rheological models, parameter identification, the operator of Riemann--Liouville, experimental data, polyvinylchloride, error}

\date{13/VII/2016}
\revisiondate{26/VIII/2016}
\accepteddate{09/IX/2016}

%%%%%%%%%%%%%%%%%%%%%%%%%%%%%%%%%%%%%%%%%%%%%%%%%%%%%%%%%%%%%%%%%%%%%%%%%%%%%%%%%%%%


\makerutitle

\Section[n]{Введение}

В работе~\cite{OgoRadUng16} на основе гипотезы В.\,Вольтерры о наследственно упругом деформируемом твердом теле и метода структурного моделирования рассмотрены дробные аналоги некоторых классических реологических моделей: модель Скотт Блэра как аналог и обобщение модели Ньютона для вязкой жидкости, а также дробные аналоги моделей Фойхта, Максвелла, Кельвина и Зенера. В указанной работе было отмечено, что если в исходном определяющем соотношении В.\,Вольтерры использовано ядро Абеля, то возникающие в определяющих соотношениях дробные производные будут являться производными Римана"--~Лиувилля на отрезке временной оси. Показана корректность классической задачи Коши для возникающих из определяющих соотношений дробных дифференциальных уравнений относительно некоторых линейных комбинаций функций напряжений и деформаций. Найдены явные решения задачи о ползучести при постоянном напряжении на стадиях нагружения и разгрузки и отмечено, что они могут быть использованы для аппроксимации экспериментальных зависимостей деформации от времени, полученных в опытах по одноосному растяжению деформируемых тел (физических сред) под действием постоянных нагрузок. Успешное решение задачи аппроксимации позволяет идентифицировать параметры реологической модели и тем самым определить вид и структуру математической модели вязкоупругого поведения изучаемой среды в форме определяющего соотношения между функциями напряжения и деформации. В работе~\cite{OgoRadUng16} авторы ограничились примером использования дробного аналога модели Фойхта для описания данных эксперимента по нагружению образцов из поливинилхлоридного пластиката, которые приведены в работе~\cite{RadchenkoGoludin}. В настоящей работе решается задача идентификации параметров математической модели вязкоупругого деформирования указанных выше  математических объектов и выполняется оценка погрешности построенных моделей.

Прежде чем перейти к описанию алгоритма построения математической модели одноосного вязкоупругого деформирования отметим, что идея замены классических вязкоупругих моделей деформируемого твердого тела наследственно упругими моделями имеет давнюю историю, восходящую к работам В.\,Вольтерры~\cite{Volterra1909, Volterra1982} и Х.\,Больцмана~\cite{Boltzman}. Она получила мощное развитие в работе Ю.\,Н.\,Работнова~\cite{Rabotnov1948}, которая по существу заложила начало школы по механике наследственно упругих твердых тел. Отметим здесь также несколько предшествующих работ~[\citen{Duffing}--\citen{Gerasimov1948}].

К настоящему моменту времени опубликовано огромное количество работ, посвященных дробным реологическим моделям (см. библ. список в~[\citen{Bulgakov}--\citen{Mainardi2010}]). Как отмечено в монографии~\cite[с.~443]{K-S-T} большинство авторов используют в определяющих соотношениях дробные производные Капуто, производные Грюнвальда"--~Летникова~\cite{Schmidt} или лиувиллевскую форму~\cite{Lewandowski} дробного интегро-дифференцирования на всей числовой оси, в то время как дробные реологические модели выражаются через дробные производные Римана"--~Лиувилля. Для нахождения явных решений дифференциальных уравнений дробного порядка, возникающих из определяющих соотношений, используется в основном преобразование Лапласа. В работе~\cite{OgoRadUng16} показано, что они могут быть найдены собственно методами дробного анализа с использованием интегральных операторов с началом в нуле или в другой точке, соответствующей, например, конечному моменту времени.

Проблема идентификации параметров дробных реологических моделей и необходимость их определения  из экспериментальных данных обсуждалась уже авторами ранних работ по наследственно упругим моделям (см.~напр.,~\cite{Gemant1936}). В монографии Ю.\,Н.\,Работнова~\cite[с.~86]{Rabotnov1977} приведен метод прямого определения параметров дробно-экспоненциального ядра ползучести, предложенный ранее авторами работы~\cite{Zvonov1968}. Однако указанный метод основан на сравнении образов преобразований Лапласа явного (аналитического) решения задачи с образом полиномиальной аппроксимации экспериментальных точек и минимизации суммы квадратов отклонений методом координатного спуска в пространстве параметров модели. В действительности проблема может быть решена непосредственно путем сопоставления экспериментальных и теоретических значений.

В многочисленных публикациях последних десяти лет проблема идентификации параметров дробных моделей в основном решается на теоретическом уровне, например, методами спектрального анализа~\cite{Vasques}. В работах~\cite{Erohin2015,Erohin_Disser} параметры модели определяются исходя из нескольких характерных точек, полученных в эксперименте, путем подстановки значений деформации в аналитические решения соответствующей задачи.

\smallskip

\Section{Описание эксперимента и идентификация параметров математических моделей}

Рассмотрим дробные аналоги известных структурных моделей вязкоупругого тела: модели Ньютона, Максвелла, Фойхта, Кельвина и Зенера с дробным порядком $\alpha\in(0,1)$. В работах ~\cite{OgoRadUng16, Abusaitova_some_special_functions} были найдены явные решения задачи о ползучести при постоянном напряжении в случае нагружения и разгрузки для всех выше упомянутых моделей. Рассмотрим применение этих математических моделей к конкретному эксперименту.
В работе~\cite{RadchenkoGoludin} исследовано поведение поливинилхлоридных трубок площадью поперечного сечения $1.2$~мм$^{2}$ при температуре 20\,\textcelsius\ и постоянных напряжениях $\sigma_{0}\in\bigl\{4.655$; $6.288$; $8.738$; $10.372$; $12.005\bigr\}$~МПа, действующих на образцы  в~течение $8$~часов, после чего производилась полная разгрузка образцов ($\sigma=0$ при $t>8$~часов).
При каждом напряжении испытывалось от 3 до 5 образцов, а затем результаты испытаний осреднялись.
Анализ экспериментальных кривых показал, что данный материал обладает следующими свойствами: имеет мгновенную упругую деформацию при приложении и снятии нагрузки, а под действием постоянной нагрузки проявляет свойство вязкоупругости. Простейшей математической моделью, описывающей подобное поведение деформируемых тел, является дробный аналог модели Максвелла
\begin{equation}\label{FracMaxwell}
\eta D^{\alpha}_{0t}\varepsilon(t)=\frac{\eta}{E_{0}}D^{\alpha}_{0t}\sigma(t)+\sigma(t),
\end{equation}
где $D^{\alpha}_{0t}\varphi=\left(d/dt\right)^{n}I^{n-\alpha}_{0t}\varphi$
~--- левосторонняя дробная производная Римана"--~Лиувилля порядка
\mbox{$\alpha>0$} функции $\varphi\left(t \right)$ ($t>0$,
$n=[\alpha]+1$, $[\alpha]$~--- целая часть числа $\alpha\in R$),
$I^{\beta}_{0t}\varphi=\left(I^{\beta}_{0+}\varphi\right)\left(t\right)$~---
интеграл Римана"--~Лиувилля порядка $\beta>0$~\cite{K-S-T}, $E_{0}$ и $\eta$"--- постоянные величины, которые при $\alpha=1$ совпадают с модулем упругого элемента и коэффициентом демпфирования соответственно в классической модели Максвелла~\cite{OgoRadUng16}.
Для напряжения, изменяющегося по закону
\begin{equation}\label{sigma0}
\sigma(t)=\sigma_{0}[H(t)-H(t-t_{1})],
\end{equation}
где $\sigma_{0}=\rm const>0$, $H(t)$~--- функция Хевисайда, $t_{1}$~--- момент времени разгрузки, решение задачи ползучести для модели~\eqref{FracMaxwell} имеет вид
\begin{equation} \label{FracMaxwellSolution}
\varepsilon(t)=\frac{\sigma_{0}}{E_{0}}\left[\left(1+\frac{\beta
t^{\alpha}}{\text Г(\alpha+1)}\right)H(t)-\Bigl(1+\frac{\beta(
t-t_{1})^{\alpha}}{\text Г(\alpha+1)}\Bigr)H(t-t_{1})\right],
\end{equation}
где $\beta=E_{0}/\eta$, $\Gamma(z)$~--- гамма-функция~\cite{Bateman}.

Рассмотрим методику идентификации параметров решения для дробного аналога модели Максвелла~\eqref{FracMaxwellSolution}. Будем аппроксимировать данные эксперимента теоретическими значениями деформации на стадии нагружения по формуле
\begin{equation} \label{FracMaxwellSolution1}
\varepsilon(t)=\frac{\sigma_{0}}{E_{0}}+\frac{\sigma_{0}t^{\alpha}}{\eta\Gamma(\alpha+1)}.
\end{equation}
Идентификации подлежат три параметра: порядок дробности $\alpha$, величины $E_{0}$ и $\eta$. Мгновенная упругая деформация выражается соотношением\\
$\varepsilon(0+)=\sigma_{0}/E_{0}$. Тогда параметр $E_{0}=\sigma_{0}/\tilde{\varepsilon}(0)$, где $\tilde{\varepsilon}(0)=\tilde{\varepsilon}_{0}$~--- первое значение деформации, взятое из массива экспериментальных данных. Для нахождения остальных неизвестных величин воспользуемся методом наименьших квадратов.

Пусть $t_{i}$"--- набор $n$ моментов времени $t$, при которых было произведено измерение экспериментальных значений деформации  $\tilde{\varepsilon}_{i}=\tilde{\varepsilon}(t_{i})$ в каждый момент времени $t_{i}$. Тогда равенство~\eqref{FracMaxwellSolution1} можно записать в виде:
\begin{equation*}
\tilde{\varepsilon}_{i}-\tilde{\varepsilon}_{0}=\frac{\sigma_{0}t^{\alpha}_{i}}{\eta\Gamma(\alpha+1)}.
\end{equation*}
Разделив левую и правую части этого равенства на величину $\sigma_{0}$ и прологарифмировав, получим следующее выражение:
\begin{equation} \label{logarifm}
\ln(\tilde{\varepsilon}_{i}-\tilde{\varepsilon}_{0})-\ln\sigma_{0}=\alpha\ln t_{i}-(\ln\eta+\ln\Gamma(\alpha+1)).
\end{equation}
Введем обозначения: $u_{i}=\ln(\tilde{\varepsilon}_{i}-\tilde{\varepsilon}_{0})-\ln\sigma_{0}$, $\tau_{i}=\ln t_{i}$, $c=\ln \eta+\ln \Gamma(\alpha+1)$. Тогда~\eqref{logarifm} запишется в виде $u_{i}=\alpha\tau_{i}-c$. Найдем параметры $\alpha$ и $c$, минимизирующие функцию
\begin{equation} \label{function2}
f=\sum^{n}_{i=1}(u_{i}-\alpha\tau_{i}+c)^{2} \rightarrow \min.
\end{equation}

Используя необходимые условия экстремума ${\partial f}/{\partial\alpha}=0$ и ${\partial f}/{\partial c}=0$ и решив систему уравнений метода наименьших квадратов, получим формулы для определения значений параметров  $\alpha$ и $c$:
$$\alpha=\frac{n\sum^{n}_{i=1}u_{i}\tau_{i}-\sum^{n}_{i=1}\tau_{i}\sum^{n}_{i=1}u_{i}}{n\sum^{n}_{i=1}\tau_{i}^{2}-\left(\sum^{n}_{i=1}\tau_{i}\right)^{2}},
$$
$$
c=\frac{\sum^{n}_{i=1}u_{i}\tau_{i}\sum^{n}_{i=1}\tau_{i}-\sum^{n}_{i=1}u_{i}\sum^{n}_{i=1}\tau_{i}^{2}}{n\sum^{n}_{i=1}\tau_{i}^{2}-\left(\sum^{n}_{i=1}\tau_{i}\right)^{2}}.
$$
Возвращаясь к введенным выше обозначениям, находим коэффициент~$\eta$:
$$
\eta={e^{c}}/{\Gamma(\alpha+1)}.
$$
Найденные значения параметров позволяют по формуле~\eqref{FracMaxwellSolution} найти теоретические значения деформации в указанные выше моменты времени.

Отметим, что если в решении задачи о ползучести в рамках дробного аналога Максвелла~\eqref{FracMaxwellSolution} отбросить мгновеннную упругую деформацию, то получим решение   соответствующей задачи для модели Скотт Блэра:
\begin{equation} \label{SB_Solution}
\varepsilon(t)=\frac{\sigma_{0}}{\eta\Gamma(\alpha+1)}[t^{\alpha}H(t)-(t-t_{1})^{\alpha}H(t-t_{1})].
\end{equation}
При аппроксимации экспериментальных данных функцией~\eqref{SB_Solution} вычтем из тестовых экспериментальных значений деформации величину мгновенной упругой деформации, совпадающую со значением $\tilde{\varepsilon}(0)$. Тогда дальнейший поиск параметров $\alpha$ и $\eta$ аналогичен алгоритму для модели Максвелла, при этом в минимизирующей функции~\eqref{function2} значения
\begin{equation}\label{u_i}
u_{i}=\ln\tilde{\varepsilon}_{i}-\ln\sigma_{0},
\end{equation}
так как $\tilde{\varepsilon}_{0}=0$. Таким образом, параметры $\alpha$ и $\eta$ в модели Скотт Блэра окажутся такими же, как и в дробном аналоге Максвелла.

Рассмотрим теперь дробный аналог модели Кельвина. В работе~\cite{OgoRadUng16} для этой модели приведено
решение в случае, когда $\sigma(t)$ определяется формулой~\eqref{sigma0}:
\begin{multline}
\label{FracKelvinSolution}
\varepsilon(t)=\frac{\sigma_{0}}{E_{2}}\Bigl[H(t)\Bigl[1+\frac{E_{2}}{E_{1}}\Bigl(1-{\rm Exp}\Bigl(\alpha,\,1;-\frac{E_{1}}{\eta};t\Bigl)\Bigr)\Bigr]-\\
-H(t-t_{1})\Bigl[1+\frac{E_{2}}{E_{1}}\Bigr(1-{\rm Exp}\Bigl(\alpha,\,1;-\frac{E_{1}}{\eta};t-t_{1}\Bigl)\Bigr)\Bigr]\Bigr],
\end{multline}
где $E_{1}$, $E_{2}$~--- коэффициенты упругих элементов, ${\rm Exp}(\alpha,\mu;\lambda;t)=\\=t^{\mu-1}E_{\alpha}(\lambda t^{\alpha};\mu)$~--- обобщенная (двупараметрическая) дробная экспоненциальная функция~\cite{Ogorodnikov_Some_Special_Functions},
 $E_{\alpha}(z;\mu)$~--- функция типа Миттаг"--~Леффлера~\cite{Djrbashyan1966}.

Из~\eqref{FracKelvinSolution} следует, что идентификации подлежат четыре коэффициента: $\alpha$, $E_{1}$, $E_{2}$, $\eta$. Вновь ограничимся случаем нагружения $\sigma(t)=\sigma_{0}=\rm const$ ($0<t\leq t_{1}$). Деформация образца описывается выражением
\begin{equation}
\label{FracKelvinSolution2}
\varepsilon(t)=\frac{\sigma_{0}}{E_{2}}+\frac{\sigma_{0}}{E_{1}}\Bigl(1-{\rm Exp}\Bigl(\alpha,\,1;-\frac{E_{1}}{\eta};t\Bigl)\Bigr).
\end{equation}
Для нахождения параметра, характеризующего мгновенную упругую деформацию, используем первую экспериментальную точку $\tilde{\varepsilon}(0)$. Тогда $E_{2}=\sigma_{0}/\tilde{\varepsilon}(0)$. Известно, что дробный аналог модели Кельвина имеет асимптоту~\cite{OgoRadUng16}
\begin{equation*}
\lim_{t\to+\infty}\varepsilon(t)={\sigma_{0}}/{E},
\end{equation*}
где $E=E_{1}E_{2}/(E_{1}+E_{2}),$ если $\sigma(t)=\sigma_{0}=\rm const$, $t\in(0,+\infty)$. Подставив в это выражение последнее экспериментальное значение деформации $\tilde{\varepsilon}_{n}$ при нагружении, определим параметр
$E_{1}\approx{\sigma_{0}}/{\tilde{\varepsilon}_{n}}$ в грубом приближении. Следующие параметры $\alpha$ и $\eta$ также в грубом приближении будут определяться с помощью метода наименьших квадратов. Для этого разложим дробную экспоненциальную функцию ${\rm Exp}(\alpha,\mu;\lambda;t)$ в ряд~\cite{Ogorodnikov_Some_Special_Functions}
\begin{multline}\label{Exp}
{\rm Exp}(\alpha,1;-\beta;t)=\sum^{\infty}_{k=0}\frac{(-\beta)^{k}t^{\alpha k}}{\Gamma(\alpha k+1)}=1+\frac{-\beta t^{\alpha}}{\Gamma(\alpha+1)}+\frac{\beta^{2}t^{2\alpha}}{\Gamma(2\alpha+1)}+\ldots,
\end{multline}
где $\beta=-{E_{1}}/{\eta}$, и ограничимся первыми двумя членами ряда~\eqref{Exp}. Тогда зависимость деформации $\varepsilon(t)$  в случае нагрузки при $0<t\leq t_{1}$ примет вид
\begin{equation}\label{KelvinToMaxwell}
\varepsilon(t)-\varepsilon(0)\approx\frac{\sigma_{0}}{E_{1}}\cdot\frac{\beta t^{\alpha}}{\Gamma(\alpha+1)}=\frac{\sigma_{0}t^{\alpha}}{\eta\Gamma(\alpha+1)}.
\end{equation}
Очевидно, что равенство~\eqref{KelvinToMaxwell} совпадает с решением задачи ползучести для дробного аналога модели Максвелла~\eqref{FracMaxwellSolution1}, для которого уже были найдены необходимые  значения параметров.
Дальнейшие уточнения параметров $\alpha$, $\eta$ и $E_{1}$ осуществляются методом покоординатного спуска (методом Хука"--~Дживса)~\cite{Huk}.
Таким образом, аппроксимация дробным аналогом модели Максвелла является грубым приближением для аппроксимации дробной моделью Кельвина.
Найденные параметры подставим в решение~\eqref{FracKelvinSolution} и вычислим среднеквадратические оценки отклонений теоретических расчетных значений от экспериментальных значений деформаций по формуле

\begin{equation}\label{norma}
\Delta=\sqrt{\displaystyle\sum^{n}_{i=1}\left[\tilde{\varepsilon}(t_{i})-\varepsilon(t_{i})\right]^{2}/\displaystyle\sum^{n}_{i=1}\left[\tilde{\varepsilon}(t_{i})\right]^{2}}\cdot 100\%
\end{equation}
на стадиях нагрузки и разгрузки.

Заметим, что исключив мгновенную упругую деформацию из уравнения~\eqref{FracKelvinSolution}, получим решение задачи о ползучести в рамках дробного аналога модели Фойхта:
\begin{multline}
\label{FracVoightSolution}
\varepsilon(t)=\frac{\sigma_{0}}{E_{1}}[H(t)(1-{\rm Exp}(\alpha,1;-\frac{E_{1}}{\eta};t))-
\\
-H(t-t_{1})(1-{\rm Exp}(\alpha,1;-\frac{E_{1}}{\eta};t-t_{1}))],
\end{multline}
Вычитая из экспериментальных данных первое значение деформации $\tilde{\varepsilon}(0)$, соответствующее упругой деформации, можно аппроксимировать полученные экспериментальные тестовые кривые уравнением~\eqref{FracVoightSolution}. Для этого следует воспользоваться идеей, изложенной выше при использовании в качестве аппроксимирующей модели дробный аналог модели Кельвина~\eqref{FracKelvinSolution}. Минимизирующая функция в этом случае совпадает с минимизирующей функцией для модели Скотт Блэра~\eqref{SB_Solution}, в которой $u_{i}$ задается формулой~\eqref{u_i}.

Рассмотрим далее дробный аналог модели Зенера, для которого решение задачи о ползучести имеет вид~\cite{OgoRadUng16}
\begin{multline}
\label{FracZenerSolution}
\varepsilon(t)=\frac{\sigma_{0}}{E_{1}+E_{2}}\Bigl[H(t)\Bigl[1+\frac{E_{1}}{E_{2}}\Bigl(1-{\rm Exp}\Bigl(\alpha,\,1;-\frac{E_{0}}{\eta};t\Bigl)\Bigr)\Bigr]-
\\
-H(t-t_{1})\Bigl[1+\frac{E_{1}}{E_{2}}\Bigr(1-{\rm Exp}\Bigl(\alpha,\,1;-\frac{E_{0}}{\eta};t-t_{1}\Bigl)\Bigr)\Bigr]\Bigr],
\end{multline}
где $E_{0}=E_{1}E_{2}/(E_{1}+E_{2})$.

Ограничимся стадией нагружения, т. е. $\sigma(t)=\sigma_{0}=\rm const$ $(0<t\leq t_{1})$. Деформация будет определяться равенством
\begin{equation} \label{FracZenerSolution2}
\varepsilon(t)=\frac{\sigma_{0}}{E_{1}+E_{2}}+\frac{E_{1}\sigma_{0}}{E_{2}(E_{1}+E_{2})}\left(1-{\rm Exp}\left(\alpha,1;-\frac{E_{0}}{\eta};t\right)\right).
\end{equation}

Выделим  мгновенную упругую компоненту деформации, благодаря которой вычисляется величина $E_{1}+E_{2}=\sigma_{0}/\tilde{\varepsilon}(0)$, где $\varepsilon(0)$~--- экспериментальное значение упругой деформации. Устремляя в уравнении~\eqref{FracZenerSolution2} $t\to+\infty$, получим $\lim_{t\to+\infty}\varepsilon(t)=\sigma_{0}/E_{2}\approx\tilde{\varepsilon}_{n}$. Откуда следует, что $E_{2}\approx\sigma_{0}/\tilde{\varepsilon}_{n}$. Тогда  $E_{1}\approx \sigma_{0}/\tilde{\varepsilon}(0)-E_{2}\approx \sigma_{0}(1/\tilde{\varepsilon}(0)-1/\tilde{\varepsilon}_{n})$.
Решение~\eqref{FracZenerSolution2} по структуре аналогично~\eqref{FracKelvinSolution2}. Вновь разложим дробную экспоненциальную функцию в ряд и ограничимся первыми двумя членами. Тогда на стадии нагрузки из~\eqref{FracZenerSolution2} получим
\begin{equation*}
\varepsilon(t)-\varepsilon(0)=\frac{\sigma_{0}E_{1}^{2}t^{\alpha}}{(E_{1}+E_{2})^{2}\eta\Gamma(\alpha+1)}.
\end{equation*}
%Откуда следует, что
%\begin{equation}
%\frac{(\varepsilon(t)-\varepsilon(0))(E_{1}+E_{2})^{2}}{\sigma_{0}E_{1}^{2}}=\frac{t^{\alpha}}{\eta\Gamma(\alpha+1)}
%\end{equation}
Минимизируя функцию~\eqref{function2}, в которой
$$
u_{i}=\ln \frac{(\varepsilon(t)-\varepsilon(0))(E_{1}+E_{2})^{2}}{\sigma_{0}E_{1}^{2}},
$$
находим коэффициенты $\alpha$ и $\eta$ в грубом приближении. Далее, следуя методу Хука"--~Дживса, аппроксимируем экспериментальные данные уравнением~\eqref{FracZenerSolution} при любом $t$ и находим отклонение теоретической зависимости $\varepsilon(t_{i})$ с найденными параметрами $\alpha$, $E_{1}$, $E_{2}$, $\eta$ от экспериментальных данных $\tilde{\varepsilon}(t_{i})$ по формуле~\eqref{norma}.

Нетрудно увидеть, что при отбрасывании начального упругого значения деформации решение задачи о ползучести для дробного аналога модели Зенера принимает вид, подобный решению для дробного аналога модели Фойхта~\eqref{FracVoightSolution}:
\begin{equation*}
%\label{FracZenerSolution2}
\varepsilon(t)=\frac{E_{1}\sigma_{0}}{E_{2}(E_{1}+E_{2})}\Bigl(1-{\rm Exp}\Bigl(\alpha,\,1;-\frac{E_{0}}{\eta};t\Bigl)\Bigr).
\end{equation*}
Последнее равенство, в свою очередь, при замене дробной экспоненты на первые два члена ее ряда принимает вид решения уравнения Скотт Блэра~\eqref{SB_Solution}, в котором значения параметров уже известны.

\smallskip

\Section{Анализ и сравнение результатов аппроксимаций}

В результате первичной обработки каждой кривой вязкоупругого деформирования при всех пяти напряжениях $\sigma_{0}$ ($\sigma(t)=\sigma_{0}=\rm const$
$(0\leq t \leq 8)$, $\sigma(t)=0$ ($8<t\leq 12$)) найдены значения параметров для всех рассмотренных моделей.
В таблице 1 приведены значения параметров аппроксимации для дробного аналога модели Кельвина. В колонке $\Delta$ с использованием нормы~\eqref{norma}  приведены отклонения теоретической зависимости $\varepsilon(t_{i})$,  рассчитанной по формуле~\eqref{FracKelvinSolution} при каждом значении напряжения $\sigma_{0}$ с найденными параметрами $\alpha$, $E_{1}$, $E_{2}$, $\eta$, от экспериментальных данных. В качестве примера на рисунке~\ref{Kelvin570g} штриховой линией представлена полученная по формуле~\eqref{FracKelvinSolution} расчетная зависимость при $\sigma_{0}=4.655$~MПa с параметрами из первой строки таблицы 1. Погрешность аппроксимации конкретной реализации в данном случае равна 2.586 \%. Анализ данных таблицы 1 свидетельствует, что, вообще говоря, для дробного аналога модели Кельвина наблюдается существенный разброс параметров модели для реализаций при различных значениях напряжений.

Аналогичный расчет проведен и для других моделей при всех пяти уровнях напряжений. Для построения исследуемых моделей, <<работоспособных>>, при напряжениях
$4.655$ МПа $\leq \sigma_{0} \leq 12.005$ МПа, выполнено усреднение параметров по пяти реализациям для каждой из них.  В таблице 2 представлены полученные значения усредненных параметров $\bar{\alpha}$, $\bar{E}_{1}$, $\bar{E}_{2}$, $\bar{\eta}$ для всех моделей.
\smallskip
\begin{figure}%[b!]
\noindent
\begin{minipage}[h]{0.48\textwidth}

\caption{Экспериментальная (сплошная линия) и расчетная по модели~(\ref{FracKelvinSolution}) (штриховая линия) кривые вязкоупругого деформирования поливинилхлоридного пластиката при напряжении $\sigma_{0}=4.655$~МПа с~последующей разгрузкой
\label{Kelvin570g}
 }

\smallskip \footnotesize

[Figure~\ref{Kelvin570g}.
Experimental (solid line) and calculated by the model~\eqref{FracKelvinSolution} (dashed line) viscoelastic creep curves of the flexible PVC at the stress $\sigma_{0}=4.655$~MPa with subsequent \centerline{unloading]}


\end{minipage} \hfill
\begin{minipage}[h]{0.5\textwidth}


\centering
\includegraphics[width=0.98\textwidth]{Kelvin570g}

\end{minipage}

\end{figure}


\smallskip

\begin{center}
 %\label{tab:KelvinError}
{\bf\footnotesize Таблица 1. Значения параметров аппроксимации~(\ref{FracKelvinSolution}) и погрешности для дробного аналога модели Кельвина
[The values \par of the approximation parameters for Eq.~(\ref{FracKelvinSolution}) and measure of inaccuracy for fractional Kelvin model] }


        \begin{tabular}{|c||c|c|c|c|c|}
\hline
\rule{0mm}{11pt}%
\footnotesize $\sigma_{0}$, МПа & \footnotesize  $\alpha$ & \footnotesize  $E_{1}$, МПа & \footnotesize  $E_{2}$, МПа & \footnotesize  $\eta$ & \footnotesize  $\Delta,\, \%$ \\
            \hline
            \rule{0mm}{12pt}%
            \phantom{1}4.655 & 0.355 & 145.525 & 1163.800 & 190.067 & 2.586 \\
             \phantom{1}6.288 & 0.326 & 133.601 & \phantom{1}898.329 & 158.793 & 4.097 \\
             \phantom{1}8.738 & 0.394 & 140.194 & 1028.000 & 128.956 & 3.168 \\
             10.372 & 0.318 & \phantom{1}97.571 & \phantom{1}829.733 & 158.840 & 3.468 \\
             12.005 & 0.400 & 127.962 & \phantom{1}857.500 & 113.856 & 3.121 \\
            \hline
        \end{tabular}

\smallskip

{\bf\footnotesize Таблица 2. Усредненные значения параметров для дробных моделей
[The averaged parameter values for the fractional models] }

\begin{tabular}{|c||c|c|c|c|} %\label{tab:ModelsParameters}
\hline
\rule{0mm}{11pt}%
\footnotesize Дробные модели & \footnotesize  $\bar{\alpha}$ & \footnotesize  $\bar{E}_{1}$, МПа & \footnotesize  $\bar{E}_{2}$, МПа & \footnotesize  $\bar{\eta}$  \\
            \hline
            \rule{0mm}{12pt}%
            Скотт Блэра & 0.179  &  --  &  --  & 289.847   \\
            Фойхта & 0.359  & \phantom{1}128.971 & -- &  150.102 \\
            Максвелла & 0.179 &  --  & 955.472 & 289.847 \\
            Кельвина & 0.359 & \phantom{1}128.971 & 955.472 & 150.102 \\
            Зенера & 0.285 & 1078.692 & \phantom{1}92.418 & 143.029 \\
            \hline
        \end{tabular}

\smallskip

{\bf\footnotesize Таблица 3. Погрешности аппроксимаций для всех исследуемых моделей после осреднения параметров для пяти реализаций
[The approximation errors for all investigated models after averaging parameters for the five realizations] }


\begin{tabular}{|c||c|c|c|c|c|} %\label{tab:ModelsError}
\hline
\rule{0mm}{11pt}%
 & \multicolumn{5}{c|}{Напряжения, МПа}\\ \hline
\footnotesize Дробные модели & \footnotesize  4.655  & \footnotesize  6.288 & \footnotesize  8.738  & \footnotesize  10.372 & \footnotesize  12.005  \\
            \hline
            \rule{0mm}{12pt}%
            Скотт Блэра & 19.139 \% & 12.538 \% & 10.192 \% & 13.237 \% & 15.745 \%  \\
            Фойхта & 18.580 \% & \phantom{1}8.892 \% & \phantom{1}4.841 \% & 11.062 \% & 12.273 \%\\
            Максвелла & 19.792 \% & \phantom{1}9.671 \% & \phantom{1}8.578 \% & 12.609 \% & 14.429 \%\\
            Кельвина & 19.325 \% & \phantom{1}6.606 \% & \phantom{1}3.964 \% & 11.153 \% & 11.843 \%\\
            Зенера & 19.018 \% & \phantom{1}7.394 \% & \phantom{1}4.937 \% & 11.438 \% & 12.527 \%\\
            \hline
        \end{tabular}

\end{center}





В таблице 3 приведены погрешности отклонений расчетных значений вязкоупругой деформации от экспериментальных для всех пяти моделей с усредненными параметрами, вычисленные по формуле~\eqref{norma} при каждом уровне напряжений. Средняя же погрешность по всем пяти реализациям для модели Скотт Блэра равна 14.170~\%, для дробных аналогов модели Фойхта "--- 11.130~\%, Максвелла "--- 13.016~\%, Кельвина "--- 10.578~\% и Зенера "--- 11.063~\%.
Таким образом, наилучшее приближение теоретической кривой к экспериментальным значениям дает дробный аналог модели Кельвина. В качестве иллюстрации на рисунках~\ref{KelvinAll} и~\ref{MaxwellAll} штриховой линией построены расчетные значения вязкоупругой деформации по моделям~\eqref{FracKelvinSolution} и~\eqref{FracMaxwellSolution} соответственно с усредненными коэффициентами.


Отметим, что все вычисления и графические построения производились в пакете MATLAB, а в расчетах использован <<Автоматизированный исследовательский комплекс
\mbox{<<MitLef>>.\footnote{Электронный ресурс ``Автоматизированный исследовательский комплекс «MitLef»''/ Н.~С.~Яшагин, Е.~Н.~Огородников,
Зарег. в ОФЭРНиО \No~17486 от 11.10.2011~г. и ФГНУ ЦИТиС \No~50201151294 от 18.11.2011~г.}}
для вычисления дробной экспоненциальной функции ${\rm Exp}(\alpha,\mu;\lambda;t)$.



\smallskip

\Section{Выводы}


Сравненительный анализ погрешностей отклонений расчетных кривых от экспериментальных для всех моделей показал, что дробный аналог модели Кельвина наилучшим образом аппроксимирует экспериментальные данные, полученные при растяжении поливинилхлоридного стержня. Однако, по результатам вычислений средней погрешности отклонений расчетных кривых с усредненными параметрами для пяти уровней напряжений по каждой модели видно, что дробная модель Кельвина с коэффициентами $\bar{\alpha}$, $\bar{E}_{1}$, $\bar{E}_{2}$, $\bar{\eta}$ не столь существенно
улучшает аппроксимацию экспериментальных данных в сравнении с другими моделями. Близкие результаты дают также дробный аналог модели Фойхта с тремя параметрами и Зенера с четырьмя параметрами.












        %\begin {tabular} {|l|l|l|l|l|} \cline{1-5}
        %\hline
        %\rule{0mm}{11pt}
        %\multicolumn{5} {|c|} {Дробные модели}\\
        %\hline
        %Скотт Блэра & Фойхта & Максвелла & Кельвина & Зенера \\
        %\hline
        %\rule{0mm}{12pt}
        %14.170 \% & 11.130 \% & 13.016 \% & 10.578 \% & 11.063 \%\\
        %\hline
        %\end{tabular}



\smallskip

\begin{figure}[p!]

\noindent
\begin{minipage}[h]{0.5\textwidth}
\centering  \scriptsize \sl
\includegraphics[width=0.98\textwidth]{Kelvin570g-a}\\
a
\end{minipage}
\begin{minipage}[h]{0.5\textwidth}
\centering  \scriptsize \sl
\includegraphics[width=0.98\textwidth]{Kelvin770g-b}\\
b
\end{minipage}

\vspace{2mm}

\noindent
\begin{minipage}[h]{0.5\textwidth}
\centering  \scriptsize \sl
\includegraphics[width=0.98\textwidth]{Kelvin1070g-c}\\
c
\end{minipage}
\begin{minipage}[h]{0.5\textwidth}
\centering  \scriptsize \sl
\includegraphics[width=0.98\textwidth]{Kelvin1270g-d}\\
d
\end{minipage}

\vspace{2mm}

\begin{minipage}[h]{0.5\textwidth}
\centering  \scriptsize \sl
\includegraphics[width=0.98\textwidth]{Kelvin1470g-e}\\
e
\end{minipage} \hfill
\begin{minipage}[h]{0.48\textwidth}
\caption{Экспериментальная (сплошная линия) и расчетная по модели~(\ref{FracKelvinSolution}) с усредненными значениями параметров
(штриховая линия) кривые вязкоупругого деформирования поливинилхлоридного пластиката при напряжениях $\sigma_{0}=4.655$~МПа ({\sl a}\/),
$\sigma_{0}=6.288$~МПа ({\sl b}\/),
$\sigma_{0}=8.738$~МПа ({\sl c}\/), $\sigma_{0}=10.372$~МПа ({\sl d}\/) и
$\sigma_{0}=12.005$~МПа ({\sl e}\/)
 с~последующей разгрузкой \label{KelvinAll} }
\end{minipage}

\smallskip \footnotesize

[Figure~\ref{KelvinAll}.
Experimental (solid line) and calculated by the model~\eqref{FracKelvinSolution}  with averaged values of parameters  (dashed line) viscoelastic creep curves of the flexible PVC at the stresses \mbox{$\sigma_{0}=4.655$~MPa} ({\sl a}\/), $\sigma_{0}=6.288$~MPa ({\sl b}\/),
$\sigma_{0}=8.738$~MPa ({\sl c}\/), $\sigma_{0}=10.372$~MPa ({\sl d}\/),  and
\centerline{$\sigma_{0}=12.005$~MPa ({\sl e}\/) with subsequent unloading]}
\end{figure}



\smallskip

        \begin{figure}[p!]

\noindent
\begin{minipage}[h]{0.5\textwidth}
\centering  \scriptsize \sl
\includegraphics[width=0.98\textwidth]{Maxwell570g-a}\\
a
\end{minipage}
\begin{minipage}[h]{0.5\textwidth}
\centering  \scriptsize \sl
\includegraphics[width=0.98\textwidth]{Maxwell770g-b}\\
b
\end{minipage}

\vspace{2mm}

\noindent
\begin{minipage}[h]{0.5\textwidth}
\centering  \scriptsize \sl
\includegraphics[width=0.98\textwidth]{Maxwell1070g-c}\\
c
\end{minipage}
\begin{minipage}[h]{0.5\textwidth}
\centering  \scriptsize \sl
\includegraphics[width=0.98\textwidth]{Maxwell1270g-d}\\
d
\end{minipage}

\vspace{2mm}

\begin{minipage}[h]{0.5\textwidth}
\centering  \scriptsize \sl
\includegraphics[width=0.98\textwidth]{Maxwell1470g-e}\\
e
\end{minipage} \hfill
\begin{minipage}[h]{0.48\textwidth}
\caption{Экспериментальная (сплошная линия) и расчетная по модели~(\ref{FracMaxwellSolution}) с усредненными значениями параметров (штриховая линия) кривые вязкоупругого деформирования поливинилхлоридного пластиката при напряжениях $\sigma_{0}=4.655$~МПа ({\sl a}\/), $\sigma_{0}=6.288$~МПа ({\sl b}\/),
$\sigma_{0}=8.738$~МПа ({\sl c}\/), $\sigma_{0}=10.372$~МПа ({\sl d}\/) и
$\sigma_{0}=12.005$~МПа ({\sl e}\/)
 с~последующей разгрузкой
\label{MaxwellAll}
 }

\end{minipage}

\smallskip \footnotesize

[Figure~\ref{MaxwellAll}.
Experimental (solid line) and calculated by the model~\eqref{FracMaxwellSolution}  with averaged values of parameters  (dashed line) viscoelastic creep curves of the flexible PVC at the stresses \mbox{$\sigma_{0}=4.655$~MPa} ({\sl a}\/), $\sigma_{0}=6.288$~MPa ({\sl b}\/),
$\sigma_{0}=8.738$~MPa ({\sl c}\/), $\sigma_{0}=10.372$~MPa ({\sl d}\/),  and
\centerline{$\sigma_{0}=12.005$~MPa ({\sl e}\/) with subsequent unloading]}
\end{figure}


\thanks{{\bf ORCIDs}

{\bf Луиза Гадильевна Унгарова:} \url{http://orcid.org/0000-0001-5388-8101}

}


\RusRef
\begin{thebibliography}{10}

\RBibitem{OgoRadUng16}
\by Огородников~Е.~Н., Радченко~В.~П., Унгарова~Л.~Г.
\paper Математическое моделирование наследственно упругого деформируемого тела
на основе структурных моделей и~аппарата дробного
интегро-дифференцирования Римана--Лиувилля
\serial Вестн. Сам. гос. техн. ун-та. Сер. Физ.-мат. науки
\yr 2016
\vol 20
\issue 1
\pages 167--194
\mathnet{http://mi.mathnet.ru/vsgtu1456}
\crossref{10.14498/vsgtu1456}
\elib{http://elibrary.ru/item.asp?id=26898215}

\RBibitem{RadchenkoGoludin}
\by Радченко~В.~П., Голудин~Е.~П.
\paper Феноменологическая стохастическая модель изотермической ползучести поливинилхлоридного пластиката
\jour Вестн. Сам. гос. техн. ун-та. Сер. Физ.-мат. науки
\yr 2008
\issue 1(16)
\pages 45--52
\mathnet{http://mi.mathnet.ru/vsgtu571}
\crossref{10.14498/vsgtu571}

\Bibitem{Volterra1909}
\paper Sulle equazioni integro-differenziali della teoria dell'elasticit\`a
\by Volterra~V.
\jour Rend. Acc. Naz. Lincei
\yr 1909
\vol 18
\pages 295–301

\RBibitem{Volterra1982}
\book Теория функционалов, интегральных и интегро-дифференциальных уравнений
\by Вольтерра~В.
\publaddr М.
\publ Наука
\yr 1982
\totalpages 304

\Bibitem{Boltzman}
\by Boltzmann~L.
\paper Theorie der elastischen Nachwirkung (Theory of elastic after effects)
\zmath{https://zbmath.org/?q=an:02717441}
\jour Wien. Ber.
\vol 70
\yr 1874
\pages  275–306\!\!
\moreref
\paper  Zur Theorie der elastischen Nachwirkung (On the elastic after effect)
\by    Boltzman~L.
\zmath{https://zbmath.org/?q=an:10.0670.01}
\jour Pogg. Ann. \rm (2)
\vol 5
\pages 430--432\!\!
\yr 1878
\moreref
\paper  Zur Theorie der elastischen Nachwirkung
\by    Boltzman~L.
\ed  Friedrich Hasen\"ohrl
\inbook Wissenschaftliche Abhandlungen
\bookvol 2
\publaddr Cambridge
\publ Cambridge University Press
\yr 2012
\serial Cambridge Library Collection
\pages 318--320
\crossref{10.1017/CBO9781139381437.015}

\RBibitem{Rabotnov1948}
\by Работнов~Ю.~Н.
\paper Равновесие упругой среды с последействием
\jour ПММ
\yr 1948
\vol 12
\issue 1
\pages 53--62
\zmath{https://zbmath.org/?q=an:03055326}

\Bibitem{Duffing}
\by Duffing~G.
\paper Elastizit\"at und Reibung beim Riementrieb (Elasticity and friction of the belt drive)
\jour Forschung auf dem Gebiet des Ingenieurwesens~A
\vol 2
\issue 3
\pages 99--104
\yr 1931
\crossref{10.1007/BF02578795}


\Bibitem{Gemant1936}
\by Gemant~A.
\paper A Method of Analyzing Experimental Results Obtained from Elasto-Viscous Bodies
\jour J. Appl. Phys.
\vol 7
\yr 1936
\pages 311--317
\crossref{10.1063/1.1745400}

\Bibitem{Gemant1938}
\by Gemant~A.
\paper On fractional differentials
\jour Philos. Mag., VII. Ser.
\vol 25
\pages 540--549
\yr 1938
\zmath{https://zbmath.org/?q=an:0018.25201}

\RBibitem{Bronsky}
\by Бронский~А.~П.
\paper Явление последействия в твердом теле
\jour ПММ
\yr 1941
\vol 5
\issue 1
\pages 31--56
\zmath{https://zbmath.org/?q=an:0063.00617}


\RBibitem{Slonimsky}
\by Слонимский~Г.~Л.
\paper О законах деформации реальных материалов
\jour ЖТФ
\yr 1939
\vol 9
\issue 20
\pages 1791--1799

\RBibitem{Gerasimov1948}
\by Герасимов~А.~Н.
\paper Обобщение линейных законов деформирования и его применение к задачам внутреннего трения
\jour  ПММ
\vol 12
\issue 3
\yr 1948
\pages 251--260
\zmath{https://zbmath.org/?q=an:0032.12901}

\RBibitem{Bulgakov}
\by Булгаков~И.~И.
\book Ползучесть полимерных материалов
\publaddr М.
\publ Наука
\yr 1973
\totalpages 288

\Bibitem{Podlubny1999}
\by Podlubny~I.
\book Fractional differential equations. An introduction to fractional derivatives, fractional differential equations, to methods of their solution and some of their applications
\zmath{https://zbmath.org/?q=an:01329842}
\serial Mathematics in Science and Engineering
\vol 198
\publaddr San Diego
\publ Academic Press
\yr 1999
\totalpages xxiv+340


\Bibitem{K-S-T}
\serial North-Holland Mathematics Studies
\vol 204
\yr 2006
\book Theory and Applications of Fractional Differential Equations
\by Kilbas~A.~A., Srivastava~H.~M., Trujillo~J.~J.
\publ Elsevier
\publaddr Amsterdam
\totalpages xx+523
\zmath{https://zbmath.org/?q=an:02217537}


\Bibitem{Mainardi2010}
\by Mainardi~F.
\book Fractional calculus and waves in linear viscoelasticity. An introduction to mathematical models
\publaddr Hackensack, NJ
\publ World Scientific
\yr 2010,
\totalpages xx+347~pp
\crossref{10.1142/9781848163300}
\zmath{https://zbmath.org/?q=an:1210.26004}

\Bibitem{Schmidt}
\by Schmidt~A., Gaul~L.
\paper Parameter Identification and FE Implementation of a Viscoelastic Constitutive Equation Using Fractional Derivatives
\jour Proc. Appl. Math. Mech.
\vol 1(1)
\yr 2002
\pages 153--154
\crossref{10.1002/1617-7061(200203)1:1<153::AID-PAMM153>3.0.CO;2-J}

\Bibitem{Lewandowski}
\by Lewandowski~R., \Russian Chor\k{a}\.zyczewski~B.
\paper  Identification of the parameters of the Kelvin--Voigt and the Maxwell fractional models, used to modeling of viscoelastic dampers
\jour Computers and Structures
\vol 88
\issue 1--2
\yr 2009
\pages 1--17
\crossref{10.1016/j.compstruc.2009.09.001}

\RBibitem{Rabotnov1977}
\book Элементы наследственной механики твердых тел
\by Работнов~Ю.~Н.
\publaddr М
\publ Наука.
\yr 1977
\totalpages 384


\RBibitem{Zvonov1968}
\by Звонов~Е.~Н., Малинин~Н.~И., Паперник~Л.~Х., Цейтлин~Б.~М.
\paper Определение характеристик ползучести линейных упруго-наследственных материалов с использованием ЭЦВМ
\jour  Изв. АН СССР, МТТ
\yr 1968
\issue 5
\pages 76–85

\Bibitem{Vasques}
\by Vasques C.M.A., Rodrigues J. Dias,  Moreira R.A.S.
\paper Experimental identification of GHM and ADF parameters for viscoelastic damping modeling
\inbook III European Conference on Computational Mechanics
\publ  Springer
\yr 2006
\crossref{10.1007/1-4020-5370-3_173}

\RBibitem{Erohin2015}
\by Ерохин~С.~В., Алероев~Т.~С., Фриштер~Л.~Ю., Колесниченко~А.~В.
\paper Параметрическая идентификация математической модели вязкоупругих
материалов с использованием производных дробного порядка
\jour International Journal for Computational Civil and Structural Engineering
\yr 2015
\issue 11(3)
\pages 82--86

\RBibitem{Erohin_Disser}
\by Ерохин~С.~В.
\book Математическое моделирование физических соотношений вязкоупругих тел с использованием методов дробного исчисления. Диссертация на соискание ученой степени кандидата технических наук
\publaddr M.
\yr 2015
\totalpages 117




\RBibitem{Abusaitova_some_special_functions}
\by Абусаитова~Л.~Г., Огородников~Е.~Н.
\paper О некоторых специальных функциях, связанных с функцией Миттаг--Леффлера, их свойствах и применении
\bookinfo Материалы X Школы молодых ученых
\inbook Нелокальные краевые задачи и проблемы современного анализа и информатики
\publaddr Нальчик
\publ  КБНЦ РАН
\yr 2012
\pages 13--15

\RBibitem{Bateman}
\by Бейтмен~Г., Эрдейи~А.
\book Высшие трансцендентные функции
\bookvol 1
\voltitle Гипергеометрическая функция, функции Лежандра
\yr 1973
\publ Наука
\publaddr M.
%\sortdate 1/1/1973
\totalpages 294

\RBibitem{Ogorodnikov_Some_Special_Functions}
\by Огородников~Е.~Н., Яшагин~Н.~С.
\paper Некоторые специальные функции в решении задачи Коши для одного дробного осцилляционного уравнения
\jour Вестн. Сам. гос. техн. ун-та. Сер. Физ.-мат. науки
\yr 2009
\issue 1(18)
\pages 276--279
\mathnet{http://mi.mathnet.ru/vsgtu685}
\crossref{10.14498/vsgtu685}

\RBibitem{Djrbashyan1966}
\by Джрбашян~М.~М.
\book Интегральные преобразования и представления функций в~комплексной области
\publaddr М.
\publ Наука
\yr 1966
\totalpages 672
\zmath{https://zbmath.org/?q=an:0154.37702}

\Bibitem{Huk}
\paper  "Direct Search" Solution of Numerical and Statistical Problems
\by Hooke~R.,  Jeeves~T.\,A.
\jour Journal of the ACM (JACM)
\issue 8(2)
\yr 1961
\pages 212--229
\crossref{10.1145/321062.321069}


\end{thebibliography}

\RuDate




\newpage

\makeentitle

\thanks{{\bf Declaration of Financial and Other Relationships.}
This work was supported by the program of the Ural Branch of RAS (project no.~15--10--1--18).
Each author has participated in the article concept development and in the manuscript writing.
The authors are absolutely responsible for submitting the final manuscript in print.
Each author has approved the final version of manuscript.
The authors have not received any fee for the article.
}

\pagebreak


\thanks{{\bf ORCIDs}

{\bf Luiza G. Ungarova:} \url{http://orcid.org/0000-0001-5388-8101}
}


\EngRef
\begin{thebibliography}{20}

\Bibitem{OgoRadUng16:eng}
\lang In Russian
\by Ogorodnikov~E.~N., Radchenko~V.~P., Ungarova~L.~G.
\paper Mathematical modeling of hereditary elastically deformable body on the basis of structural models and fractional integro-differentiation Riemann-Liouville apparatus
\jour Vestn. Samar. Gos. Tekhn. Univ. Ser. Fiz.-Mat. Nauki \rm [J. Samara State Tech. Univ., Ser. Phys. \& Math. Sci.]
\yr 2016
\vol 20
\issue 1
\pages 167--194
\mathnet{http://mi.mathnet.ru/vsgtu1456}
\crossref{10.14498/vsgtu1456}
\elib{http://elibrary.ru/item.asp?id=26898215}

\RBibitem{RadchenkoGoludin:eng}
\lang In Russian
\by Radchenko~V.~P., Goludin~E.~P.
\paper Phenomenological stochastic isothermal creep model for an polivinylchloride elastron
\jour Vestn. Samar. Gos. Tekhn. Univ. Ser. Fiz.-Mat. Nauki  \rm [J. Samara State Tech. Univ., Ser. Phys. \& Math. Sci.]
\yr 2008
\issue 1(16)
\pages 45--52
\crossref{10.14498/vsgtu571}

\Bibitem{Volterra1909:eng}
\paper Sulle equazioni integro-differenziali della teoria dell'elasticit\`a
\by Volterra~V.
\jour Rend. Acc. Naz. Lincei
\yr 1909
\vol 18
\pages 295–301
\zmath{https://zbmath.org/?q=an:40.0870.01}

\Bibitem{Volterra1982:eng}
\by Volterra~V.
\book Theory of functionals and of integral and integro-differential equations
\publaddr New York
\publ Dover Publ., Inc.
\totalpages 226
\yr 1959
\zmath{https://zbmath.org/?q=an:0086.10402|02573174}

\Bibitem{Boltzman}
\by Boltzmann~L.
\paper Theorie der elastischen Nachwirkung (Theory of elastic after effects)
\zmath{https://zbmath.org/?q=an:02717441}
\jour Wien. Ber.
\vol 70
\yr 1874
\pages  275–306\!\!
\moreref
\paper  Zur Theorie der elastischen Nachwirkung (On the elastic after effect)
\by    Boltzman~L.
\zmath{https://zbmath.org/?q=an:10.0670.01}
\jour Pogg. Ann. \rm (2)
\vol 5
\pages 430--432\!\!
\yr 1878
\moreref
\paper  Zur Theorie der elastischen Nachwirkung
\by    Boltzman~L.
\ed  Friedrich Hasen\"ohrl
\inbook Wissenschaftliche Abhandlungen
\bookvol 2
\publaddr Cambridge
\publ Cambridge University Press
\yr 2012
\serial Cambridge Library Collection
\pages 318--320
\crossref{10.1017/CBO9781139381437.015}

\Bibitem{Rabotnov1948:eng}
\jour Fractional Calculus and Applied Analysis
\yr 2014
\vol 17
\issue 3
\pages 684--696
\paper  Equilibrium of an elastic medium with after-effect
\by Rabotnov~Yu.~N.
\crossref{10.2478/s13540-014-0193-1}
\zmath{https://zbmath.org/?q=an:1306.74011}


\Bibitem{Duffing}
\by Duffing~G.
\paper Elastizit\"at und Reibung beim Riementrieb (Elasticity and friction of the belt drive)
\jour Forschung auf dem Gebiet des Ingenieurwesens~A
\vol 2
\issue 3
\pages 99--104
\yr 1931
\crossref{10.1007/BF02578795}


\Bibitem{Gemant1936}
\by Gemant~A.
\paper A Method of Analyzing Experimental Results Obtained from Elasto-Viscous Bodies
\jour J. Appl. Phys.
\vol 7
\yr 1936
\pages 311--317
\crossref{10.1063/1.1745400}

\Bibitem{Gemant1938}
\by Gemant~A.
\paper On fractional differentials
\jour Philos. Mag., VII. Ser.
\vol 25
\pages 540--549
\yr 1938
\zmath{https://zbmath.org/?q=an:0018.25201}

\Bibitem{Bronsky:eng}
\lang In Russian
\by Bronskij~A.~P.
\paper Residual effect in rigid bodies
\jour Prikl. Mat. Mekh.
\yr 1941
\vol 5
\issue 1
\pages 31--56


\Bibitem{Slonimsky:eng}
\by Slonimsky~G.~L.
\paper On the laws of deformation of real materials. I. The theories of Maxwell and Boltzmann
\jour Acta physicochim. URSS
\vol 12
\pages 99--128
\yr 1940
\zmath{https://zbmath.org/?q=an:66.1023.04}

\Bibitem{Gerasimov1948:eng}
\lang In Russian
\by Gerasimov A.~N.
\paper A generalization of linear laws of deformation and its application to internal friction problem
\jour  Prikl. Mat. Mekh.
\vol 12
\issue 3
\yr 1948
\pages 251--260
\zmath{https://zbmath.org/?q=an:0032.12901}

\Bibitem{Bulgakov:eng}
\by Bulgakov~I.~I.
\book Polzuchest' polimernykh materialov \rm [Creep of Polymer Materials]
\publaddr Moscow
\publ Nauka
\yr 1973
\totalpages 288

\Bibitem{Podlubny1999}
\by Podlubny~I.
\book Fractional differential equations. An introduction to fractional derivatives, fractional differential equations, to methods of their solution and some of their applications
\zmath{https://zbmath.org/?q=an:01329842}
\serial Mathematics in Science and Engineering
\vol 198
\publaddr San Diego
\publ Academic Press
\yr 1999
\totalpages xxiv+340


\Bibitem{K-S-T}
\serial North-Holland Mathematics Studies
\vol 204
\yr 2006
\book Theory and Applications of Fractional Differential Equations
\by Kilbas~A.~A., Srivastava~H.~M., Trujillo~J.~J.
\publ Elsevier
\publaddr Amsterdam
\totalpages xx+523
\zmath{https://zbmath.org/?q=an:02217537}


\Bibitem{Mainardi2010}
\by Mainardi~F.
\book Fractional calculus and waves in linear viscoelasticity. An introduction to mathematical models
\publaddr Hackensack, NJ
\publ World Scientific
\yr 2010,
\totalpages xx+347~pp
\crossref{10.1142/9781848163300}
\zmath{https://zbmath.org/?q=an:1210.26004}

\Bibitem{Schmidt}
\by Schmidt~A., Gaul~L.
\paper Parameter Identification and FE Implementation of a Viscoelastic Constitutive Equation Using Fractional Derivatives
\jour Proc. Appl. Math. Mech.
\vol 1(1)
\yr 2002
\pages 153--154
\crossref{10.1002/1617-7061(200203)1:1<153::AID-PAMM153>3.0.CO;2-J}

\Bibitem{Lewandowski}
\by Lewandowski~R., \Russian Chor\k{a}\.zyczewski~B.
\paper  Identification of the parameters of the Kelvin--Voigt and the Maxwell fractional models, used to modeling of viscoelastic dampers
\jour Computers and Structures
\vol 88
\issue 1--2
\yr 2009
\pages 1--17
\crossref{10.1016/j.compstruc.2009.09.001}

\Bibitem{Rabotnov1977:eng}
\zmath{https://zbmath.org/?q=an:0515.73026}
\book Elements of hereditary solid mechanics
\publaddr
\publ Mir Publ.
\totalpages  388
\yr 1980


\RBibitem{Zvonov1968:eng}
\lang In Russian
\by Zvonov~E.~N., Malinin~N.~I., Papernik~L.~H., Ceytlin~B.~M.
\paper Determination of the creep characteristics of linear elastic-hereditary materials using an electronic digital computer
\jour  Izv. AN SSSR, MTT
\yr 1968
\issue 5
\pages 76–85

\Bibitem{Vasques}
\by Vasques C.M.A., Rodrigues J. Dias,  Moreira R.A.S.
\paper Experimental identification of GHM and ADF parameters for viscoelastic damping modeling
\inbook III European Conference on Computational Mechanics
\publ  Springer
\yr 2006
\crossref{10.1007/1-4020-5370-3_173}

\RBibitem{Erohin2015:eng}
\lang In Russian
\by Erokhin~S.~V., Aleroev~T.~S., Frishter~L.~Y., Kolesnichenko~A.~V.
\paper Parametric identification of mathematical models of viscoelastic
materials using fractional derivatives
\jour International Journal for Computational Civil and Structural Engineering
\yr 2015
\issue 11(3)
\pages 82--86

\RBibitem{Erohin_Disser:eng}
\lang In Russian
\by Erokhin~S.~V.
\book Mathematical modeling of physical relationships of viscoelastic bodies using fractional calculus methods. Thesis on the degree for candidate of technical sciences
\publaddr M.
\yr 2015
\totalpages 117



\Bibitem{Abusaitova_some_special_functions:eng}
\lang In Russian
\by Abusaitova~L.~G., Ogorodnikov~E.~N.
\paper Some special functions associated with the Mittag--Leffler function, their properties, and applications
\inbook Nelokal'nye kraevye zadachi i problemy sovremennogo analiza i informatiki \rm [Non-local boundary value problems and problems of modern analysis and informatics]
\publaddr Nal'chik
\yr 2012
\pages 13--15


\RBibitem{Bateman}
\by Bateman~G., Erdelyi ~A.
\book Higher Transcendental Functions
\bookvol 1
\yr 1953
\publ McGraw-Hill
\publaddr New York
\totalpages 302

\Bibitem{Ogorodnikov_Some_Special_Functions:eng}
\lang In Russian
\by Ogorodnikov~E.~N., Yashagin~N.~S.
\paper Some Special Functions in the Solution To Cauchy Problem for a Fractional Oscillating Equation
\jour Vestn. Samar. Gos. Tekhn. Univ. Ser. Fiz.-Mat. Nauki \rm [J. Samara State Tech. Univ., Ser. Phys. \& Math. Sci.]
\yr 2009
\issue 1(18)
\pages 276--279
\crossref{10.14498/vsgtu685}

\Bibitem{Djrbashyan1966:eng}
\lang In Russian
\by Dzhrbashyan~M.~M.
\book Integral'nye preobrazovaniia i predstavleniia funktsii v~kompleksnoi oblasti \rm [Integral Transforms and Representation of Functions in Complex Domain]
\publaddr Moscow
\publ Nauka
\yr 1966
\totalpages 672

\Bibitem{Huk}
\paper  "Direct Search" Solution of Numerical and Statistical Problems
\by Hooke~R.,  Jeeves~T.\,A.
\jour Journal of the ACM (JACM)
\issue 8(2)
\yr 1961
\pages 212--229
\crossref{10.1145/321062.321069}

\end{thebibliography}

\EngDate


\end{document}
